\documentclass{article}
\usepackage[T1]{fontenc}
\usepackage{lmodern}
\usepackage{graphicx}
\usepackage{amsmath,amssymb}
\usepackage[a4paper,margin=2cm]{geometry}
\usepackage[absolute,overlay]{textpos}
\setlength{\TPHorizModule}{1mm}
\setlength{\TPVertModule}{1mm}
\usepackage[indonesian]{babel}
\usepackage{hyperref}
\setlength{\emergencystretch}{2em}
% unit: Halaman 16

\begin{document}

% === Halaman 16 (python/native) ===

16

Iterasi 1:


UZ QSO VUOHXMOPV GPOZPEVSG ZWSZ OPFPESX UDBMETSX AIZ


 t            e   e te     t  t  e e               t 


VUEPHZ HMDZSHZO WSFP APPD TSVP QUZW YMXUZUHSX


   e t    t  t     e  ee     e   t      t



EPYEPOPDZSZUFPO MB ZWP FUPZ HMDJ UD TMOHMQ


  e  e e t t  e     t e   et  

•ZWP dan ZWSZ dipetakan menjadi t*e dan t**t 

•Kemungkinan besar W adalah pemetataan dari H sehingga kata yang 

mungkin untuk ZWP dan ZWSZ adalah the dan that

\end{document}
